\section{Bottom normal form and the fast Lyapunov oracle}
\label{sec:normal-form}

The unique equilibrium in \cref{prop:regular-shells} is elliptic.  Its exact
linear structure both selects the orbit used in the trace and supplies
analytic limits against which the validated computation can be audited.

\subsection{Singular values, frequencies, and periods}

At the origin,
\begin{equation}
 A_a=D\Psi_a(0)=
 \begin{pmatrix}-c_a&-1\\1&0\end{pmatrix},
 \qquad
 D^2V_a(0)=4\pi^2A_a^TA_a.
 \label{eq:bottom-hessian}
\end{equation}
Let \(0<s_-<s_+\) be the singular values of \(A_a\).  Since
\(\det A_a=1\), direct diagonalization gives
\begin{equation}
 s_+s_-=1,
 \qquad
 s_\pm=\frac{\sqrt{c_a^2+4}\pm|c_a|}{2}.
 \label{eq:singular-values}
\end{equation}
In orthonormal configuration coordinates \((Q_-,Q_+)\) aligned with the
eigenvectors of \(A_a^TA_a\), the quadratic Hamiltonian is
\begin{equation}
 K_0=\frac12\left(P_-^2+\omega_-^2Q_-^2
                      +P_+^2+\omega_+^2Q_+^2\right),
 \qquad \omega_\pm=2\pi s_\pm.
 \label{eq:harmonic-normal-form}
\end{equation}
We use the signs to label frequency, so the fast mode is \(+\).  Its period
and the slow period are
\begin{equation}
 T_+^0=s_+^{-1}=s_-,
 \qquad
 T_-^0=s_-^{-1}=s_+.
 \label{eq:linear-periods}
\end{equation}
Set
\begin{equation}
 \rho_a=\frac{\omega_+}{\omega_-}=\frac{s_+}{s_-}=s_+^2>1.
 \label{eq:frequency-ratio}
\end{equation}

\begin{proposition}[Fast Lyapunov branch]
\label{prop:lyapunov-branch}
For every \(a>-1\), \(a\ne0\), a periodic family
\(\gamma_+(E)\) emanates from the fast normal mode as \(E\downarrow2\pi\).
It can be parameterized by every sufficiently small positive energy excess,
and
\begin{equation}
 T_+(E)\longrightarrow T_+^0.
 \label{eq:period-limit}
\end{equation}
The reduced transverse multipliers tend to
\(e^{\pm2\pi i/\rho_a}\), so the family is transversally nondegenerate for
all sufficiently small positive excesses.  Moreover,
\begin{align}
 D_+^0
 &:=\lim_{E\downarrow2\pi}|\det(I-P_+(E))|
 =4\sin^2\frac{\pi}{\rho_a},
 \label{eq:det-limit}\\
 S_+(E)
 &:=\oint_{\gamma_+(E)}p\dd q
 =T_+^0(E-2\pi)+o(E-2\pi).
 \label{eq:action-limit}
\end{align}
\end{proposition}

\begin{proof}
The linearized Hamiltonian vector field has the two simple pairs
\(\pm i\omega_-\) and \(\pm i\omega_+\).  For the fast pair, the other
frequency ratios are
\(\pm\omega_-/\omega_+=\pm1/\rho_a\notin\mathbb Z\).  The
Lyapunov-centre theorem therefore produces a one-parameter family
\citep{AlligoodYorke1986,Weinstein1973}.  If \(A\) is the fast
normal-coordinate amplitude, then
\(E(A)-2\pi=\tfrac12\omega_+^2A^2+O(A^3)>0\) for small \(A>0\), so energy
indexes the positive branch.  The multiplier and determinant limits follow
from the slow rotation during one fast period.  Finally, the standard action
identity \(\dd S_+/\dd E=T_+\), with \(S_+(2\pi)=0\), gives
\cref{eq:action-limit}.
\end{proof}

\subsection{First nonlinear period coefficient}

The leading harmonic limits are insufficient to audit a finite positive
energy.  Let
\begin{equation}
 C_{ijk}=\partial_{Q_iQ_jQ_k}V_a(0),
 \qquad
 D_{ijkl}=\partial_{Q_iQ_jQ_kQ_l}V_a(0)
 \label{eq:tensors}
\end{equation}
in the fixed orthonormal normal-coordinate convention.  Define, for
\(j\in\{-,+\}\),
\begin{equation}
 b_{j0}=-\frac{C_{++j}}{4\omega_j^2},
 \qquad
 b_{j2}=-\frac{C_{++j}}{4(\omega_j^2-4\omega_+^2)},
 \label{eq:b-coefficients}
\end{equation}
and
\begin{equation}
 \nu_+=\frac{1}{2\omega_+}
 \left[
 \frac{D_{++++}}{8}
 +\sum_{j\in\{-,+\}}C_{++j}
 \left(b_{j0}+\frac12b_{j2}\right)
 \right].
 \label{eq:nu-plus}
\end{equation}

\begin{proposition}[Nonlinear period slope]
\label{prop:period-slope}
With the tensors and coordinates above,
\begin{equation}
 \left.\frac{\dd T_+}{\dd E}\right|_{2\pi+}
 =-\frac{2T_+^0\nu_+}{\omega_+^3}.
 \label{eq:period-slope}
\end{equation}
For \(a=1.02\), the deterministic coordinate convention gives
\(\nu_+=17.52709598189346\) and hence
\begin{align}
 T_+(2\pi+\delta)
 &=0.6638439766792985
   -0.0274450756283701\,\delta+o(\delta),
 \label{eq:period-expansion}\\
 \frac{S_+(2\pi+\delta)}{\delta}
 &=0.6638439766792985
   -0.0137225378141851\,\delta+o(\delta).
 \label{eq:action-expansion}
\end{align}
\end{proposition}

The proof is a Poincar\'e--Lindstedt solvability calculation.  Its key point
is that the constant and second harmonics induced at order \(A^2\) feed the
resonant \(A^3\cos\tau\) equation.  Omitting that feedback changes the
period slope.  We give the calculation in \cref{app:nonlinear-period}.

\subsection{Numerical specialization at \texorpdfstring{\(a=1.02\)}{a=1.02}}

The exact formulas give
\begin{equation}
 T_+^0=0.6638439766792985,
 \qquad
 T_-^0=1.5063780573896775,
 \qquad
 D_+^0=3.8627220445155035.
 \label{eq:oracle-numbers}
\end{equation}
In particular,
\begin{equation}
 T_+^0<0.75<\min(2T_+^0,T_-^0).
 \label{eq:window-separation}
\end{equation}
This strict ordering is the limiting return classification behind the
whole-shell theorem in \cref{sec:analytic-trace}.  The limiting geometric
amplitude is
\begin{equation}
 \frac{T_+^0}{\sqrt{D_+^0}}=0.3377686126427769,
 \qquad
 \frac{T_+^0}{2\pi\sqrt{D_+^0}}=0.0537575443233896,
 \label{eq:limiting-amplitudes}
\end{equation}
where the second value matches the Fourier normalization in
\cref{eq:fourier-convention}.

A nonvalidated continuation run at six positive excesses served as a
formula and convention audit.  The extrapolated values and analytic oracles
are shown in \cref{tab:r400-oracle}.  These comparisons do not prove the
branch or the trace theorem; those statements come from the analytic and
validated arguments below.

\begin{table}[t]
\centering
\caption{Small-energy continuation audit.  ``Fitted'' means a deterministic
extrapolation of the archived nonvalidated orbit data; the analytic column
is independent of the fit.}
\label{tab:r400-oracle}
\small
\begin{tabular}{@{}lrrr@{}}
\toprule
Quantity & Fitted & Analytic & Absolute difference\\
\midrule
\(T_+^0\) & 0.663843973386761 & 0.663843976679299 & \(3.29\times10^{-9}\)\\
\(\dd T_+/\dd E\) & -0.0274445154485 & -0.0274450756284 & \(5.60\times10^{-7}\)\\
\(\lim S/\delta\) & 0.663843975854219 & 0.663843976679299 & \(8.25\times10^{-10}\)\\
\(\dd(S/\delta)/\dd\delta\) & -0.0137223974763 & -0.0137225378142 & \(1.40\times10^{-7}\)\\
\(D_+^0\) & 3.86272204305148 & 3.86272204451550 & \(1.46\times10^{-9}\)\\
\bottomrule
\end{tabular}
\end{table}
